%% kapa: exam=1 round=26 year=2015 period=2 form=A subject=accounting subject_ko=회계학 seq=4 number=44
%% source: Q-Net 감정평가사 기출문제 (cst00309) artlSeq=5138001 file=extracted/1008/5138001/제26회_감정평가사 1차_기출/제26회_감정평가사 1차_기출/2교시_감평_A형.pdf
%% conversion: kordoc 4.24.51 -> tools/gen1.py (auto); review=proofread-sonnet; flags=none
\begin{question}{44}
(주)감평은 20×1년 12월 말 현재 창고에 보관 중인 상품재고를 실사한 결과, 금액이 ￦2,000,000임을 확인하였다. 다음 자료를 반영하여 계산한 기말상품재고액은?
\begin{framedtext}
○ 거래처에 FOB 도착지 인도조건으로 발송한 상품(원가 ￦400,000)이 20×1년 말 현재 운송 중에 있다. \par
○ 20×1년 12월 28일에 고객에게 발송한 시송품(원가 ￦200,000) 가운데 20×1년 말까지 매입의사표시를 해 온 금액이 ￦120,000이다. \par
○ 20×1년 12월 15일에 거래처에 반품가능판매 조건으로 상품(원가 ￦500,000)을 판매하였다. 12월 말까지 반품된 상품은 없으며, 반품률은 10\%로 신뢰성 있게 추정된다.
\end{framedtext}
\begin{choices}
  \item ￦2,130,000
  \item ￦2,480,000
  \item ￦2,530,000
  \item ￦2,980,000
  \item ￦3,100,000
\end{choices}
\end{question}
\answer{2}
